\documentclass[border=10pt]{standalone}
\usepackage{tikz,ctex}
\usetikzlibrary{positioning}

\begin{document}

\begin{tikzpicture}[
    % 全局设置
    node distance = 1.6cm and 1.4cm,
    every node/.style = {
        draw,
        rounded corners,
        thick,
        align = center,
        font = \sffamily\small,
        minimum width = 2.0cm,
        minimum height = 0.9cm
    },
    % 样式定义
    centerStyle/.style = {
        fill = blue!45,
        text = white,
        font = \sffamily\bfseries\large,
        line width = 2pt,
        minimum size = 2.5cm,
        circle
    },
    branchTitle/.style = {
        font = \sffamily\bfseries,
        line width = 1.5pt
    },
    subNode/.style = {
        fill = white,
        font = \sffamily\small,
        minimum width = 2.0cm,
        minimum height = 0.8cm
    },
    leafNode/.style = {
        fill = white,
        font = \sffamily\footnotesize,
        minimum width = 1.8cm,
        minimum height = 0.7cm
    },
    arrowStyle/.style = {
        ->,
        thick,
        >=stealth,
        shorten >=2pt
    }
]

    % ================= 1. 中心节点 =================
    \node[centerStyle] (center) {自动微分 \\ Automatic Differentiation};

    % ================= 2. 上方分支：四种求导方法 (红) =================
    \node[branchTitle, fill=red!15, above=of center] (methods) {四种求导方法 \\ Four Approaches};

    \node[subNode, above left=of methods] (manual) {手动反向传播 \\ Manual Backprop};
    \node[subNode, above=of methods] (numeric) {数值微分 \\ Numerical Diff};
    \node[subNode, above right=of methods] (symbolic) {符号微分 \\ Symbolic Diff};
    \node[subNode, right=of methods, yshift=0cm] (autodiff) {自动微分 \\ Autodiff};

    % ================= 3. 右侧分支：前向模式 (蓝) =================
    \node[branchTitle, fill=orange!15, right=of center] (forward) {前向模式 \\ Forward Mode};

    \node[subNode, above=of forward] (primal) {主变量 $v_i$ \\ Primal Variables};
    \node[subNode, above right=of forward] (tangent) {切线变量 $\dot{v}_i$ \\ Tangent Variables};
    \node[subNode, right=of forward] (trace) {求值轨迹 \\ Evaluation Trace};
    \node[subNode, below right=of forward, yshift=0cm] (jvp) {雅可比--向量积 \\ JVP};
    \node[subNode, below=of forward] (column) {雅可比一列 \\ One Jacobian Column};

    % ================= 4. 下方分支：反向模式 (绿) =================
    \node[branchTitle, fill=green!15, below=of center] (reverse) {反向模式 \\ Reverse Mode};

    \node[subNode, below left=of reverse] (adjoint) {伴随变量 $\bar{v}_i$ \\ Adjoint Variables};
    \node[subNode, below=of reverse] (backward) {反向传播 \\ Backward Pass};
    \node[subNode, below right=of reverse] (row) {雅可比一行 \\ One Jacobian Row};
    \node[subNode, left=of reverse, yshift=0cm] (memory) {内存开销大 \\ High Memory};
    \node[subNode, right=of reverse] (vjp) {向量--雅可比积 \\ VJP};

    % ================= 5. 左侧分支：对比与应用 (紫) =================
    \node[branchTitle, fill=purple!15, left=of center] (compare) {对比与应用 \\ Comparison \& Applications};

    \node[subNode, above left=of compare] (complexity) {复杂度对比 \\ Complexity};
    \node[subNode, left=of compare] (hybrid) {混合模式 \\ Hybrid Mode};
    \node[subNode, below left=of compare] (hessian) {Hessian--向量积 \\ Hessian-Vector};
    \node[subNode, above=of compare, yshift=-0.3cm] (cost) {常数倍开销 \\ Constant Overhead};

    % ================= 连线逻辑 =================
    % 中心 -> 四大分支标题
    \draw[arrowStyle, red] (center) -- (methods);
    \draw[arrowStyle, blue] (center) -- (forward);
    \draw[arrowStyle, green!60!black] (center) -- (reverse);
    \draw[arrowStyle, purple] (center) -- (compare);

    % 分支标题 -> 子节点 (上方)
    \draw[arrowStyle, red!60] (methods) -- (manual);
    \draw[arrowStyle, red!60] (methods) -- (numeric);
    \draw[arrowStyle, red!60] (methods) -- (symbolic);
    \draw[arrowStyle, red!60] (methods) -- (autodiff);

    % 分支标题 -> 子节点 (右侧)
    \draw[arrowStyle, blue!60] (forward) -- (primal);
    \draw[arrowStyle, blue!60] (forward) -- (tangent);
    \draw[arrowStyle, blue!60] (forward) -- (trace);
    \draw[arrowStyle, blue!60] (forward) -- (jvp);
    \draw[arrowStyle, blue!60] (forward) -- (column);

    % 分支标题 -> 子节点 (下方)
    \draw[arrowStyle, green!60!black] (reverse) -- (adjoint);
    \draw[arrowStyle, green!60!black] (reverse) -- (backward);
    \draw[arrowStyle, green!60!black] (reverse) -- (row);
    \draw[arrowStyle, green!60!black] (reverse) -- (memory);
    \draw[arrowStyle, green!60!black] (reverse) -- (vjp);

    % 分支标题 -> 子节点 (左侧)
    \draw[arrowStyle, purple!60] (compare) -- (complexity);
    \draw[arrowStyle, purple!60] (compare) -- (hybrid);
    \draw[arrowStyle, purple!60] (compare) -- (hessian);
    \draw[arrowStyle, purple!60] (compare) -- (cost);

\end{tikzpicture}

\end{document}